\section{Introduction}

% Move 1: state the complete problem.
During real-world robot deployment, a policy produces an action trajectory from
the current observation, and executing that trajectory changes the physical
scene. Collecting an alternative action under exactly the same conditions would
require reproducing object poses, contact states, sensor readings, and timing.
Such reconstruction is generally unavailable after execution. The deployed
trajectory can nevertheless be inspected afterward and labeled as desirable or
undesirable. The resulting deployment feedback contains one executed trajectory
and one unary judgment for each encountered context, but no same-context
counterfactual. It identifies behavior to encourage or suppress without
providing the alternative needed for a direct comparison. Reusing this feedback
is attractive precisely because it avoids recollecting counterfactual actions
or rewarded candidate groups on real hardware. The challenge is to turn this
direction-only deployment feedback into a valid policy update.

Flow-based robot policies add a separate optimization constraint. They generate
continuous action trajectories by integrating a learned vector field from noise
through multiple steps \cite{lipman2023flowmatching,chi2023diffusionpolicy}.
Their networks therefore do not emit terminal action likelihoods as direct
outputs. Naively differentiating a trajectory-level objective through the full
generation process can also yield unstable gradients when adapting such a
deployed policy. The problem is consequently defined by both the deployment
feedback and the policy interface: can a pretrained flow-based robot
policy be improved directly from unary feedback collected during real-world
deployment, without collecting additional robot interactions for optimization
or learning a reward or value critic?

% Move 2: review existing methods and identify the mismatch.
Existing post-training methods obtain a learning signal by adding comparison or
value structure. Direct Preference Optimization requires a preferred and a
dispreferred output for the same context \cite{rafailov2023dpo}; Proximal Policy
Optimization and Group Relative Policy Optimization instead form advantages
from rewarded current-policy trajectories or groups of candidates
\cite{schulman2017ppo,schulman2016gae,shao2024deepseekmath}. Recent flow-policy
methods illustrate critic-based alternatives. RECAP converts deployment
outcomes into binarized advantage conditioning through a learned value function,
using demonstrations, autonomous rollouts, and human corrections
\cite{physicalintelligence2025pi06star}. Q-learning with Adjoint Matching learns
a temporal-difference critic from reward-labeled transitions and uses adjoint
matching to turn its action gradients into supervision for the flow field
\cite{li2026qam}. These routes are viable under their respective data
interfaces, but they do not define the restricted update studied here, whose
only quality signal is a fixed set of individually judged deployment
trajectories. They first construct pairwise or group comparisons, collect
current-policy experience, or learn critic-derived targets. We instead ask
whether unary judgments can reach the flow policy without these intermediaries.

% Move 3: state the proposed response.
Without a learned intermediary, a unary judgment must be converted into a
comparison using quantities available from the policy itself. The first
comparison asks how the current policy's fit to a labeled deployment trajectory
has changed relative to the frozen pretrained policy. This reveals how training
has moved the recorded behavior, but not how that change stands relative to the
current policy's own outputs. A second, policy-centered comparison must therefore
come from trajectories generated by the current policy at the stored deployment
contexts. Recorded deployment actions cannot provide this reference because
they represent the behavior-data distribution; actions borrowed from other
contexts would additionally confound policy change with context--action
mismatch. Together, the two comparisons anchor each update to the pretrained
policy while interpreting it against the policy being optimized.

We propose Flow-KTO to implement this principle. Flow-KTO first measures the
change in conditional flow-matching fit for each labeled deployment trajectory
relative to a frozen reference policy. It then samples trajectories internally
from the current policy at stored deployment contexts, evaluates each sample
under the current and reference policies, and aggregates those comparisons into
a dynamic reference point. This construction keeps every labeled comparison
anchored to the pretrained policy while locating it within the current policy's
output distribution. The unary label determines which side of this
policy-centered reference should receive utility, and the exact bounded KTO
objective converts that direction into a policy update
\cite{ethayarajh2024kto}. The internally sampled trajectories are not assigned
quality labels and are never executed on the robot; they are used only to
estimate the current policy's reference point. Because
flow-matching differences need not share the numerical scale of explicit
likelihood ratios, Flow-KTO calibrates the inverse temperature to the scale of
these comparisons while keeping class-utility and objective-mixing weights
separate.

% Add the controlled-results summary only after the Experiments section fixes
% the matched baselines, estimator ablations, uncertainty, and robot protocol.
% Move 4: itemize the core contributions.
The paper makes three technical contributions:

\begin{itemize}
    \item We formulate direct post-training of a pretrained flow-based robot
    policy from a fixed collection of feedback gathered during real-world robot
    deployment, with one unary-labeled trajectory per encountered context and no
    requirement for paired alternatives, rewarded groups, additional robot
    interaction, or a learned reward or value critic.
    \item We introduce a policy-centered reference estimator in the native
    flow-matching training space. It uses internally sampled current-policy
    trajectories at stored deployment contexts and is distinct from references
    computed on recorded deployment actions or context-mismatched actions.
    \item We optimize the resulting label-directed comparison with exact
    bounded KTO and calibrate its inverse temperature to the scale of the
    flow-matching comparisons, separate from class-utility and objective-mixing
    weights.
\end{itemize}
